% --- LaTeX Lecture Notes Template - S. Venkatraman ---

% --- Set document class and font size ---

\documentclass[letterpaper, 10pt]{article}


% --- Package imports ---

% Extended set of colors
\usepackage[dvipsnames]{xcolor}

\usepackage{
  amsmath, amsthm, amssymb, mathtools, dsfont, units,          % Math typesetting
  graphicx, wrapfig, subfig, float,                            % Figures and graphics formatting
  listings, color, inconsolata, pythonhighlight,               % Code formatting
  fancyhdr, sectsty, hyperref, enumerate, enumitem, framed }   % Headers/footers, section fonts, links, lists

% lipsum is just for generating placeholder text and can be removed
\usepackage{hyperref, lipsum} 
% \usepackage{bm}
% \usepackage{subfig}

% --- Fonts ---

\usepackage{newpxtext, newpxmath, inconsolata}
\usepackage{amsfonts}
\usepackage{pgfplots}
\pgfplotsset{compat=1.12}
\usepackage{tkz-fct}
\usepackage{svg}
\usepackage{tikz}
\usepackage{tikz-cd}
\usepackage{lipsum}
\usepackage{enumitem}
\usepackage[title]{appendix}
% \usepackage[toc,page]{appxendix}
\usepackage[utf8]{inputenc}
\usepackage{multicol}
\usepackage{multirow}
\usepackage{booktabs}

\usepackage{minted}  % code highlighting
% \usepackage[finalizecache,cachedir=.]{minted}
% \usepackage[frozencache,cachedir=.]{minted}

\DeclareUnicodeCharacter{3BC}{$\pi$}
\DeclareUnicodeCharacter{3C0}{$\pi$}

\usepackage[most]{tcolorbox}
\newtcolorbox{myquote}[1][]{%
  colback=black!5,
  colframe=black!5,
  notitle,
  sharp corners,
  % borderline west={2pt}{0pt}{red!80!black},
  enhanced,
  breakable,
}

\renewcommand*\pod[1]{%
  \allowbreak
  \mathchoice
    {\mkern 18mu}%
    {\mkern  8mu}%
    {\mkern  6mu}% "6mu" matches the space *after* the word "mod"
    {\mkern  6mu}%
  (#1)%
}

% --- Page layout settings ---

% Set page margins
\usepackage[left=1.35in, right=1.35in, top=1.0in, bottom=.9in, headsep=.2in, footskip=0.35in]{geometry}

% Anchor footnotes to the bottom of the page
\usepackage[bottom]{footmisc}

% Set line spacing
\renewcommand{\baselinestretch}{1.2}

% Set spacing between paragraphs
\setlength{\parskip}{1.3mm}

% Allow multi-line equations to break onto the next page
\allowdisplaybreaks

% --- Page formatting settings ---

% Set image captions to be italicized
\usepackage[font={it,footnotesize}]{caption}

% Set link colors for labeled items (blue), citations (red), URLs (orange)
\hypersetup{colorlinks=true, linkcolor=RoyalBlue, citecolor=RedOrange, urlcolor=ForestGreen}

% Set font size for section titles (\large) and subtitles (\normalsize) 
\usepackage{titlesec}
% \titleformat{\section}{\large\bfseries}{{\fontsize{19}{19}\selectfont\textreferencemark}\;\; }{0em}{}
\titleformat{\section}{\large\bfseries}{\thesection\;\;\;}{0em}{}
\titleformat{\subsection}{\normalsize\bfseries\selectfont}{\thesubsection\;\;\;}{0em}{}

% Enumerated/bulleted lists: make numbers/bullets flush left
%\setlist[enumerate]{wide=2pt, leftmargin=16pt, labelwidth=0pt}
\setlist[itemize]{wide=0pt, leftmargin=16pt, labelwidth=10pt, align=left}

% --- Table of contents settings ---

\usepackage[subfigure]{tocloft}

% Reduce spacing between sections in table of contents
\setlength{\cftbeforesecskip}{.9ex}

% Remove indentation for sections
\cftsetindents{section}{0em}{0em}

% Set font size (\large) for table of contents title
\renewcommand{\cfttoctitlefont}{\large\bfseries}

% Remove numbers/bullets from section titles in table of contents
\makeatletter
\renewcommand{\cftsecpresnum}{\begin{lrbox}{\@tempboxa}}
\renewcommand{\cftsecaftersnum}{\end{lrbox}}
\makeatother

% --- Set path for images ---

\graphicspath{{Images/}{../Images/}}

% --- Math/Statistics commands ---

% Add a reference number to a single line of a multi-line equation
% Usage: "\numberthis\label{labelNameHere}" in an align or gather environment
\newcommand\numberthis{\addtocounter{equation}{1}\tag{\theequation}}

% Shortcut for bold text in math mode, e.g. $\b{X}$
\let\b\mathbf

% Shortcut for bold Greek letters, e.g. $\bg{\beta}$
\let\bg\boldsymbol

% Shortcut for calligraphic script, e.g. %\mc{M}$
\let\mc\mathcal

% \mathscr{(letter here)} is sometimes used to denote vector spaces
\usepackage[mathscr]{euscript}

% Convergence: right arrow with optional text on top
% E.g. $\converge[p]$ for converges in probability
\newcommand{\converge}[1][]{\xrightarrow{#1}}

% Weak convergence: harpoon symbol with optional text on top
% E.g. $\wconverge[n\to\infty]$
\newcommand{\wconverge}[1][]{\stackrel{#1}{\rightharpoonup}}

% Equality: equals sign with optional text on top
% E.g. $X \equals[d] Y$ for equality in distribution
\newcommand{\equals}[1][]{\stackrel{\smash{#1}}{=}}

% Normal distribution: arguments are the mean and variance
% E.g. $\normal{\mu}{\sigma}$
\newcommand{\normal}[2]{\mathcal{N}\left(#1,#2\right)}

% Uniform distribution: arguments are the left and right endpoints
% E.g. $\unif{0}{1}$
\newcommand{\unif}[2]{\text{Uniform}(#1,#2)}

% Independent and identically distributed random variables
% E.g. $ X_1,...,X_n \iid \normal{0}{1}$
\newcommand{\iid}{\stackrel{\smash{\text{iid}}}{\sim}}

% Sequences (this shortcut is mostly to reduce finger strain for small hands)
% E.g. to write $\{A_n\}_{n\geq 1}$, do $\bk{A_n}{n\geq 1}$
\newcommand{\bk}[2]{\{#1\}_{#2}}

\newcommand{\what}[1]{\widehat{#1}}
% \setcounter{section}{-1}

\setcounter{MaxMatrixCols}{20}

\newcommand{\SL}{\mathrm{SL}}
\newcommand{\Sp}{\mathrm{Sp}}
\newcommand{\Mp}{\mathrm{Mp}}
\newcommand{\GL}{\mathrm{GL}}
\newcommand{\SO}{\mathrm{SO}}
\newcommand{\SU}{\mathrm{SU}}
\newcommand{\PGL}{\mathrm{PGL}}
\newcommand{\PSL}{\mathrm{PSL}}
\newcommand{\GSp}{\mathrm{GSp}}
\newcommand{\PGSp}{\mathrm{PGSp}}
\newcommand{\Spin}{\mathrm{Spin}}
\newcommand{\gl}{\mathfrak{gl}}
\newcommand{\JL}{\mathrm{JL}}
\newcommand{\stab}{\mathrm{Stab}}
% \newcommand{\ab}{\mathrm{ab}}
\newcommand{\cha}{\mathrm{char}}
\newcommand{\fin}{\mathrm{fin}}
\newcommand{\Tr}{\mathrm{Tr}}
\newcommand{\Li}{\mathrm{Li}}
\newcommand{\trdeg}{\mathrm{trdeg}}
\newcommand{\rank}{\mathrm{rank}}
\newcommand{\rad}{\mathrm{rad}}
\newcommand{\Res}{\mathrm{Res}}
\newcommand{\Hom}{\mathrm{Hom}}
\newcommand{\Spec}{\mathrm{Spec}\,}
\newcommand{\Sym}{\mathrm{Sym}}
\newcommand{\ev}{\mathrm{ev}}
\newcommand{\disc}{\mathrm{disc}}
\newcommand{\Aut}{\mathrm{Aut}}
\newcommand{\Span}{\mathrm{Span}}
\newcommand{\supp}{\mathrm{supp}}
\newcommand{\sgn}{\mathrm{sgn}}
\newcommand{\Lie}{\mathrm{Lie}}
\newcommand{\Ind}{\mathrm{Ind}}
\newcommand{\pInd}{\mathrm{pInd}}
\newcommand{\Gal}{\mathrm{Gal}}
\newcommand{\Cl}{\mathrm{Cl}}
\newcommand{\Wh}{\mathrm{Wh}}
\newcommand{\std}{\mathrm{std}}
\newcommand{\Slit}{\mathrm{Slit}}
\newcommand{\pprod}{\sideset{}{'}\prod}
\newcommand{\potimes}{\sideset{}{'}\otimes}
\newcommand{\pbigotimes}{\sideset{}{'}\bigotimes}
% \sideset{}{'}\prod
\newcommand{\RQM}{\mathcal{RQM}}

\newcommand{\QM}{\mathcal{QM}}

\newcommand{\dd}{\mathrm{d}}
\newcommand{\dso}{\mathds{1}}

\newcommand{\llb}{\llbracket}
\newcommand{\rrb}{\rrbracket}

\newcommand{\rA}{\mathrm{A}}
\newcommand{\rB}{\mathrm{B}}
\newcommand{\rC}{\mathrm{C}}
\newcommand{\rD}{\mathrm{D}}
\newcommand{\rE}{\mathrm{E}}
\newcommand{\rF}{\mathrm{F}}
\newcommand{\rG}{\mathrm{G}}
\newcommand{\rH}{\mathrm{H}}
\newcommand{\rI}{\mathrm{I}}
\newcommand{\rJ}{\mathrm{J}}
\newcommand{\rK}{\mathrm{K}}
\newcommand{\rL}{\mathrm{L}}
\newcommand{\rM}{\mathrm{M}}
\newcommand{\rN}{\mathrm{N}}
\newcommand{\rO}{\mathrm{O}}
\newcommand{\rP}{\mathrm{P}}
\newcommand{\rQ}{\mathrm{Q}}
\newcommand{\rR}{\mathrm{R}}
\newcommand{\rS}{\mathrm{S}}
\newcommand{\rT}{\mathrm{T}}
\newcommand{\rU}{\mathrm{U}}
\newcommand{\rV}{\mathrm{V}}
\newcommand{\rW}{\mathrm{W}}
\newcommand{\rX}{\mathrm{X}}
\newcommand{\rY}{\mathrm{Y}}
\newcommand{\rZ}{\mathrm{Z}}

\newcommand{\bA}{\mathbb{A}}
\newcommand{\bB}{\mathbb{B}}
\newcommand{\bC}{\mathbb{C}}
\newcommand{\bD}{\mathbb{D}}
\newcommand{\bE}{\mathbb{E}}
\newcommand{\bF}{\mathbb{F}}
\newcommand{\bG}{\mathbb{G}}
\newcommand{\bH}{\mathbb{H}}
\newcommand{\bI}{\mathbb{I}}
\newcommand{\bJ}{\mathbb{J}}
\newcommand{\bK}{\mathbb{K}}
\newcommand{\bL}{\mathbb{L}}
\newcommand{\bM}{\mathbb{M}}
\newcommand{\bN}{\mathbb{N}}
\newcommand{\bO}{\mathbb{O}}
\newcommand{\bP}{\mathbb{P}}
\newcommand{\bQ}{\mathbb{Q}}
\newcommand{\bR}{\mathbb{R}}
\newcommand{\bS}{\mathbb{S}}
\newcommand{\bT}{\mathbb{T}}
\newcommand{\bU}{\mathbb{U}}
\newcommand{\bV}{\mathbb{V}}
\newcommand{\bW}{\mathbb{W}}
\newcommand{\bX}{\mathbb{X}}
\newcommand{\bY}{\mathbb{Y}}
\newcommand{\bZ}{\mathbb{Z}}

\newcommand{\bx}{\mathbf{x}}
\newcommand{\by}{\mathbf{y}}

\newcommand{\cA}{\mathcal{A}}
\newcommand{\cB}{\mathcal{B}}
\newcommand{\cC}{\mathcal{C}}
\newcommand{\cD}{\mathcal{D}}
\newcommand{\cE}{\mathcal{E}}
\newcommand{\cF}{\mathcal{F}}
\newcommand{\cG}{\mathcal{G}}
\newcommand{\cH}{\mathcal{H}}
\newcommand{\cI}{\mathcal{I}}
\newcommand{\cJ}{\mathcal{J}}
\newcommand{\cK}{\mathcal{K}}
\newcommand{\cL}{\mathcal{L}}
\newcommand{\cM}{\mathcal{M}}
\newcommand{\cN}{\mathcal{N}}
\newcommand{\cO}{\mathcal{O}}
\newcommand{\cP}{\mathcal{P}}
\newcommand{\cQ}{\mathcal{Q}}
\newcommand{\cR}{\mathcal{R}}
\newcommand{\cS}{\mathcal{S}}
\newcommand{\cT}{\mathcal{T}}
\newcommand{\cU}{\mathcal{U}}
\newcommand{\cV}{\mathcal{V}}
\newcommand{\cW}{\mathcal{W}}
\newcommand{\cX}{\mathcal{X}}
\newcommand{\cY}{\mathcal{Y}}
\newcommand{\cZ}{\mathcal{Z}}

\newcommand{\scA}{\mathscr{A}}
\newcommand{\scB}{\mathscr{B}}
\newcommand{\scC}{\mathscr{C}}
\newcommand{\scD}{\mathscr{D}}
\newcommand{\scE}{\mathscr{E}}
\newcommand{\scF}{\mathscr{F}}
\newcommand{\scG}{\mathscr{G}}
\newcommand{\scH}{\mathscr{H}}
\newcommand{\scI}{\mathscr{I}}
\newcommand{\scJ}{\mathscr{J}}
\newcommand{\scK}{\mathscr{K}}
\newcommand{\scL}{\mathscr{L}}
\newcommand{\scM}{\mathscr{M}}
\newcommand{\scN}{\mathscr{N}}
\newcommand{\scO}{\mathscr{O}}
\newcommand{\scP}{\mathscr{P}}
\newcommand{\scQ}{\mathscr{Q}}
\newcommand{\scR}{\mathscr{R}}
\newcommand{\scS}{\mathscr{S}}
\newcommand{\scT}{\mathscr{T}}
\newcommand{\scU}{\mathscr{U}}
\newcommand{\scV}{\mathscr{V}}
\newcommand{\scW}{\mathscr{W}}
\newcommand{\scX}{\mathscr{X}}
\newcommand{\scY}{\mathscr{Y}}
\newcommand{\scZ}{\mathscr{Z}}

\newcommand{\frA}{\mathfrak{A}}
\newcommand{\frB}{\mathfrak{B}}
\newcommand{\frC}{\mathfrak{C}}
\newcommand{\frD}{\mathfrak{D}}
\newcommand{\frE}{\mathfrak{E}}
\newcommand{\frF}{\mathfrak{F}}
\newcommand{\frG}{\mathfrak{G}}
\newcommand{\frH}{\mathfrak{H}}
\newcommand{\frI}{\mathfrak{I}}
\newcommand{\frJ}{\mathfrak{J}}
\newcommand{\frK}{\mathfrak{K}}
\newcommand{\frL}{\mathfrak{L}}
\newcommand{\frM}{\mathfrak{M}}
\newcommand{\frN}{\mathfrak{N}}
\newcommand{\frO}{\mathfrak{O}}
\newcommand{\frP}{\mathfrak{P}}
\newcommand{\frQ}{\mathfrak{Q}}
\newcommand{\frR}{\mathfrak{R}}
\newcommand{\frS}{\mathfrak{S}}
\newcommand{\frT}{\mathfrak{T}}
\newcommand{\frU}{\mathfrak{U}}
\newcommand{\frV}{\mathfrak{V}}
\newcommand{\frW}{\mathfrak{W}}
\newcommand{\frX}{\mathfrak{X}}
\newcommand{\frY}{\mathfrak{Y}}
\newcommand{\frZ}{\mathfrak{Z}}

\newcommand{\fra}{\mathfrak{a}}
\newcommand{\frb}{\mathfrak{b}}
\newcommand{\frc}{\mathfrak{c}}
\newcommand{\frd}{\mathfrak{d}}
\newcommand{\fre}{\mathfrak{e}}
\newcommand{\frf}{\mathfrak{f}}
\newcommand{\frg}{\mathfrak{g}}
\newcommand{\frh}{\mathfrak{h}}
\newcommand{\fri}{\mathfrak{i}}
\newcommand{\frj}{\mathfrak{j}}
\newcommand{\frk}{\mathfrak{k}}
\newcommand{\frl}{\mathfrak{l}}
\newcommand{\frm}{\mathfrak{m}}
\newcommand{\frn}{\mathfrak{n}}
\newcommand{\fro}{\mathfrak{o}}
\newcommand{\frp}{\mathfrak{p}}
\newcommand{\frq}{\mathfrak{q}}
\newcommand{\frr}{\mathfrak{r}}
\newcommand{\frs}{\mathfrak{s}}
\newcommand{\frt}{\mathfrak{t}}
\newcommand{\fru}{\mathfrak{u}}
\newcommand{\frv}{\mathfrak{v}}
\newcommand{\frw}{\mathfrak{w}}
\newcommand{\frx}{\mathfrak{x}}
\newcommand{\fry}{\mathfrak{y}}
\newcommand{\frz}{\mathfrak{z}}

\newcommand{\sage}{\raisebox{-0.16\height}{\includegraphics[height=1em]{./sagelogo.png}}\,\,}

% Math mode symbols for common sets and spaces. Example usage: $\R$
\newcommand{\R}{\mathbb{R}}	% Real numbers
\newcommand{\C}{\mathbb{C}}	% Complex numbers
\newcommand{\Q}{\mathbb{Q}}	% Rational numbers
\newcommand{\Z}{\mathbb{Z}}	% Integers
\newcommand{\N}{\mathbb{N}}	% Natural numbers
\newcommand{\F}{\mathcal{F}}	% Calligraphic F for a sigma algebra
\newcommand{\El}{\mathcal{L}}	% Calligraphic L, e.g. for L^p spaces

% Math mode symbols for probability
\newcommand{\pr}{\mathbb{P}}	% Probability measure
\newcommand{\E}{\mathbb{E}}	% Expectation, e.g. $\E(X)$
\newcommand{\var}{\text{Var}}	% Variance, e.g. $\var(X)$
\newcommand{\cov}{\text{Cov}}	% Covariance, e.g. $\cov(X,Y)$
\newcommand{\corr}{\text{Corr}}	% Correlation, e.g. $\corr(X,Y)$
\newcommand{\B}{\mathcal{B}}	% Borel sigma-algebra

% Other miscellaneous symbols
\newcommand{\tth}{\text{th}}	% Non-italicized 'th', e.g. $n^\tth$
\newcommand{\Oh}{\mathcal{O}}	% Big-O notation, e.g. $\O(n)$
\newcommand{\1}{\mathds{1}}	% Indicator function, e.g. $\1_A$

\newcommand{\nul}{\mathrm{null}}
\newcommand{\range}{\mathrm{range}}

% Additional commands for math mode
\newcommand*{\argmax}{argmax}		% Argmax, e.g. $\argmax_{x\in[0,1]} f(x)$
\newcommand*{\argmin}{argmin}		% Argmin, e.g. $\argmin_{x\in[0,1]} f(x)$
\newcommand*{\spann}{Span}		% Span, e.g. $\spann\{X_1,...,X_n\}$
\newcommand*{\bias}{Bias}		% Bias, e.g. $\bias(\hat\theta)$
\newcommand*{\ran}{ran}			% Range of an operator, e.g. $\ran(T) 
\newcommand*{\dv}{d\!}			% Non-italicized 'with respect to', e.g. $\int f(x) \dv x$
\newcommand*{\diag}{diag}		% Diagonal of a matrix, e.g. $\diag(M)$
\newcommand*{\trace}{Tr}		% Trace of a matrix, e.g. $\trace(M)$
% \newcommand*{\supp}{supp}		% Support of a function, e.g., $\supp(f)$

% Numbered theorem, lemma, etc. settings - e.g., a definition, lemma, and theorem appearing in that 
% order in Lecture 2 will be numbered Definition 2.1, Lemma 2.2, Theorem 2.3. 
% Example usage: \begin{theorem}[Name of theorem] Theorem statement \end{theorem}
\theoremstyle{definition}
\newtheorem{theorem}{Theorem}[section]
\newtheorem{proposition}[theorem]{Proposition}
\newtheorem{lemma}[theorem]{Lemma}
\newtheorem{corollary}[theorem]{Corollary}
\newtheorem{definition}[theorem]{Definition}
\newtheorem{example}[theorem]{Example}
\newtheorem{question}[theorem]{Question}
\newtheorem{conjecture}[theorem]{Conjecture}
\newtheorem{exercise}[subsubsection]{Exercise}

\theoremstyle{remark}
\newtheorem{remark}[theorem]{Remark}

% Un-numbered theorem, lemma, etc. settings
% Example usage: \begin{lemma*}[Name of lemma] Lemma statement \end{lemma*}
\newtheorem*{theorem*}{Theorem}
\newtheorem*{proposition*}{Proposition}
\newtheorem*{lemma*}{Lemma}
\newtheorem*{corollary*}{Corollary}
\newtheorem*{definition*}{Definition}
\newtheorem*{example*}{Example}
\newtheorem*{remark*}{Remark}
\newtheorem*{claim}{Claim}
\newtheorem*{conjecture*}{Conjecture}

% --- Left/right header text (to appear on every page) ---

% Do not include a line under header or above footer
\pagestyle{fancy}
\renewcommand{\footrulewidth}{0pt}
\renewcommand{\headrulewidth}{0pt}

% Right header text: Lecture number and title
\renewcommand{\sectionmark}[1]{\markright{#1} }
% \fancyhead[R]{\small\textit{\nouppercase{\rightmark}}}

% Left header text: Short course title, hyperlinked to table of contents
% \fancyhead[L]{\hyperref[sec:contents]{\small Matrix multplication}}

% \pgfkeys{/Dynkin diagram,
% edge length=1.5cm,
% fold radius=1cm,
% indefinite edge/.style={
% draw=black,
% fill=white,
% thin,
% densely dashed}}

\usepackage{fancyhdr}
\pagestyle{fancy}
\fancyhead[R]{}
\fancyhead[L]{}

% --- Document starts here ---

\begin{document}

% --- Main title and subtitle ---

\title{How modular forms fly?}
% \\[1em]
% \normalsize Re-typed by Seewoo Lee\footnote{seewoo5@berkeley.edu}}

% --- Author and date of last update ---

\author{Seewoo Lee}
\date{\normalsize\vspace{-1ex} Last updated: \today}
% \date{}

% --- Add title and table of contents ---

\maketitle

% --- Abstracts ---

% \tableofcontents\label{sec:contents}
\begin{abstract}
    This is an expository note on \emph{murmurations}, which was initially discovered by He, Lee, Oliver, and Pozdnyakov for elliptic curves \cite{he2024murmurations}.
    We focus on the cases where the mumuration density is computed, starting from the work of Zubrilina \cite{zubrilina2023murmurations}.
\end{abstract}

% --- Main content: import lectures as subfiles ---
\section{Introduction}
\label{sec:intro}

The goal of this note is to introduce the arithmetic of function fields, which is the analogue of number theory for polynomials.
Especially, our main goal is to study various evidences of the following claim:

\begin{myquote}
A theorem that holds for integers is also true for polynomials (over finite fields), and latter is often easier to prove.
\end{myquote}
For example, we will see a proof of Fermat's Last Theorem for polynomials, which only requires few pages to prove.

Dictionary between the integers and the polynomials over finite fields can be found in Table \ref{tab:dictionary} of Appendix.

\subsection*{Prerequisites}
We assume that the readers are familiar with undergraduate level algebra (groups, rings, fields, etc.), number theory (congruences, prime numbers, etc.), and a bit of complex analysis.
Some of the theory of finite fields will be reviewed in Appendix \ref{subsec:handbook_galois_ff}.


\subsection*{Notations}

Let $p$ be a prime number. We denote by $\bF_p$ the finite field of order $p$, which is the field with $p$ elements.
We denote the polynomial ring $\bF_p[T]$ by $A$.
For each nonzero polynomial $f \in A$, we denote it's norm by $|f| = p^{\deg (f)}$, where $\deg (f)$ is the degree of $f$, and we set $|0| = 0$.

\subsection*{Sage codes}

There are some codes in this note, which are mostly written in \href{https://www.sagemath.org/}{Sage}.
Sage is a free \href{https://github.com/sagemath/sage}{open-source} mathematics software system, which is built on top of many existing open-source packages and wrappedn in a Python interface.
You can run them online in \href{https://sagecell.sagemath.org/}{SageMathCell}, or install it on your computer.
Especially, a lot of number-theoretic functions are implemented in Sage, so it is much easier to experiment with it than writing your own code from scratch.
For example, to check if a large number is prime, you can simply run
\begin{minted}[fontsize=\footnotesize,framesep=2mm,bgcolor=lightgray!20]{python}
is_prime(10 ^ 9 + 7)
\end{minted}
Several useful Sage functions are listed in Appendix \ref{subsec:handbook_sage}.


\subsection*{Acknowledgements}


\begin{exercise}
    Prove that $\bZ$ is not a polynomial ring over a field. In other words, show that there is no field $k$ such that $\bZ \cong k[T]$ as rings.
\end{exercise}

\begin{exercise}
    Think about your favorite theorems in number theory, and try to find their polynomial analogues. Some of them may appear in this note, but some of them may not.
\end{exercise}


\newpage
\section{Murmuration of modular forms}

In this section, we sketch Zubrilina's computation of murmuration density for modular forms.

\subsection{Statement}

Before we state the result, we define some notations first.
\begin{itemize}
    \item For $n \in \bZ_{\ge 0}$, \emph{Chebyshev polynomial of the second kind} is defined as
    \[
        U_n(\cos \theta) = \frac{\sin((n+1) \theta)}{\sin\theta}. 
    \]
    \item For $r \in \bZ_{\ge 1}$, define
    \[
        \nu(r) := \prod_{p | r} \left(1 + \frac{p^2}{p^4 - 2p^2 - p + 1}\right)
    \]
    \item Define constants $\alpha, \beta, \gamma$ as
    \begin{align*}
        \alpha &:= 2\pi \prod_{p} \frac{p^4 - 2p^2 - p + 1}{p^4 - 2p^2 + p}, \\
        \beta &:= 2\pi \prod_{p} \frac{p^3 + p^2 - 1}{p(p^2 + p - 1)}, \\
        \gamma &:= 12 \prod_{p} \frac{p(p + 1)}{p^2 + p - 1}.
    \end{align*}
\end{itemize}
\begin{theorem}[Zubrilina {\cite{zubrilina2023murmurations}}]
    \label{thm:zubrilina_modform}
    Let $X, Y, P$ be parameters going infinite with $X, Y > 0$ and $P$ prime; assume further that $Y = (1 + o(1))X^{1 - \delta_2}$ and $P \ll X^{1 + \delta_1}$ for some $\delta_1, \delta_2$ with $2\delta_1 < \delta_2 < 1$.
    Let $y = P/X$. Then
    \begin{equation}
        \label{eqn:zubrilina_modform}
        \frac{\sum_{N \in [X, X + Y]}^\square \sum_{f \in H^{\new}(N, k)} \sqrt{P} \lambda_f(P) \varepsilon(f)}{\sum_{N \in  [X, X + Y]}^\square \sum_{f \in H^{\new}(N, k)} 1} = \cM_k(y) + O_\varepsilon \left(X^{-\delta' + \varepsilon} + \frac{1}{P}\right)
    \end{equation}
    where
    \begin{equation}
        \cM_k(y) = \frac{\alpha(-1)^{k/2 - 1}}{k - 1} \sum_{1 \le r \le 2\sqrt{y}} \nu(r) \sqrt{4y - r^2} U_{k - 2} \left(\frac{r}{2\sqrt{y}}\right) + \frac{\beta}{k-1} \sqrt{y} - \gamma \delta_{k=2} y.
    \end{equation}
\end{theorem}

\subsection{Eichler--Selberg trace formula}

To prove Theorem \ref{thm:zubrilina_modform}, one need to understand how to estimate the numerator on the LHS.
Recall that $a_f(P) = P^{(k-1)/2} \lambda_f(P)$ is the $P$-th Fourier coefficient of $f$, which is also the eigenvalue of the Hecke operator $T_P$ acting on $f$.
Also, $(-1)^{k/2}\varepsilon(f)$ is equal to the eigenvalue of the Atkin--Lehner involution $W_N = T_N$ acting on $f$.
Thus the sum appears in the numerator of LHS of \eqref{eqn:zubrilina_modform} can be interpreted as the trace of the operator $(-1)^{k/2} T_P \circ W_N$ acting on the space of cusp forms of weight $k$ and level $N$ (multiplied by $P^{1 - k/2}$).
Eichler \cite{eichler1955class} studied such a sum of traces and proved that it can be expressed in terms of (Hurwitz) class numbers, which is generalized by Selberg \cite{selberg1956harmonic}.
To account the root number $\varepsilon(f)$, i.e. eigenvalue of $W_N$, we used the following version of Eichler--Selberg trace formula by Skoruppa and Zagier \cite{skoruppa1987jacobi}.

\begin{theorem}[Skoruppa--Zagier \cite{skoruppa1987jacobi}]
    \label{thm:sz_tf}
    For square-free $N$ and prime $P \nmid N$,
    \begin{align*}
        &\sum_{f \in H^\new(N, k)} \sqrt{P} \lambda_f(P) \varepsilon(f)\nonumber \\
        &= \frac{H_1(-4PN)}{2} + (-1)^{k/2 - 1} U_{k - 2} \left(\frac{r\sqrt{N}}{2\sqrt{P}}\right) \sum_{0 < r \le 2\sqrt{P/N}} H_1(r^2 N^2 - 4PN) \\
        &\quad - \delta_{k-2}(P+1)
    \end{align*}
    Here $H_1(-d)$ ($d > 0$) is the Hurwitz class number, the number of equivalence classes of positive definite binary quadratic forms of discriminant $-d$ weighted by the number of automorphisms, i.e. with forms correspond to $x^2 + y^2$ or $x^2 + xy + y^2$ counted with multiplicity $1/2$ and $1/3$ respectively.
\end{theorem}

Hurwitz class number can be expressed as a sum of usual class numbers as
\[
H_1(-d) = \sum_{f^2 | d} h(-d / f^2) + O(1)
\]
where the ``error term'' $O(1)$ disappears if $d \ne 3 \cdot \square$ or $4 \cdot \square$.
Using this, we can rewrite the Skoruppa--Zagier trace formula as
\begin{align*}
    &\sum_{f \in H^\new(k, N)} \sqrt{P} \lambda_f(P) \varepsilon(f) \\
    &= \frac{h(-4PN)}{2} + \frac{h(-PN)}{2} - \delta_{k=2} P + O(1) \\
    &+ (-1)^{k/2 - 1} U_{k-2} \left(\frac{r\sqrt{N}}{2\sqrt{P}}\right) \sum_{1 \le r \le 2\sqrt{P/N}} \sum_{d^2 | r^2 N - 4P} h\left(\frac{N(r^2 N - 4P)}{d^2}\right)
\end{align*}

From this, our new goal is to estimage the average of class numbers over short intervals, i.e. when $N \in [X, X + Y]$ with $Y = o(X)$.
The main idea is to use class number formula to write class numbers as special $L$-values at $s = 1$, e.g.
\[
h(-d) = \frac{\sqrt{d}}{\pi} L(1, \chi_d)
\]
when $d > 4$ and $-d \equiv 1 \Mod{4}$.
Then the sum (average) of the corresponding $L$-values can be estimated via truncation and Polya--Vinogradov inequality.
For example, we have an estimate
\[
L(1, \chi_d) = \sum_{n \ge 1} \frac{\chi_d(n)}{n} = \sum_{1 \le n \le T} \frac{\chi_d(n)}{n} + O\left(\frac{\sqrt{d} \log d}{T}\right).
\]
With some hard analysis, one get the following estimations.
\begin{proposition}[Zubrilina {\cite[Proposition 3.1]{zubrilina2023murmurations}}]
    Let $P$ be an odd prime and let $[X, X + Y]$ be an interval with $Y = o(X)$.
    Then as $X \to \infty$,
    \begin{align*}
        &\frac{\zeta(2) \pi}{XY} \sideset{}{{}^\square}\sum_{\substack{N \in [X, X + Y] \\ P \nmid N}}  \left(\frac{h(-PN)}{2} + \frac{h(-4PN)}{2} \right) \\
        &= A \sqrt{y} + O_{\varepsilon} \left(\frac{1}{P^{3/2}X^{1/2}} + \frac{P^{7/12}}{Y^{5/6} X^{5/12}} + \frac{YP^{1/2}}{X^{3/2}}\right) (XYP)^{\varepsilon}
    \end{align*}
\end{proposition}


\begin{proposition}[Zubrilina {\cite[Proposition 3.2]{zubrilina2023murmurations}}]
    Let $P$ be an odd prime, $r \in \bN$, and $X > Y > 0$ be such that $r^2(X + Y) < 4 P$ for each $r > 2 \sqrt{P/X}$. Let $y = P/X$. Then
    \begin{align*}
        &\frac{\zeta(2) \pi}{ XY} \sum_{r \le 2 \sqrt{P/X}} \,\,\,\, \sideset{}{{}^\square}\sum_{\substack{N \in [X, X + Y] \\ P \nmid N}}  H_1\left(r^2 N^2 - 4 PN\right) \\
        &= \sum_{r \le 2 \sqrt{P/X}} \nu(r) \sqrt{4y - r^2} \\
        &+ O\left(\frac{P^{11/10}}{Y^{2/5}X^{9/10}} + \frac{YP}{X^2} + \frac{PY^{1/2}}{X^{3/2}} + \frac{P}{X^{1/2} Y^{13/18}} + \frac{P}{XY^{1/9}}\right) (XYP)^\varepsilon
    \end{align*}
\end{proposition}
\section{Other known cases}

After the success of Zubrilina, a lot of people are interested in murmuration density for different objects in number theory.
We list the known works here.

\subsection{Flying Dirichlet characters}

Lee, Oliver, and Pozdnyakov computed murmuration density for Dirichlet characters \cite{lee2025murmurations}\footnote{which can be thought as  automorphic forms on $\GL_1$ over $\bQ$.}.
For complex characters, the corresponding murmuration densities are given by the following theorem.

\begin{theorem}[{Lee--Oliver--Pozdnyakov \cite[Theorem 1.1]{lee2025murmurations}}]
    Let $\cD_{+}(N)$ (resp. $\cD_{-}(N)$) denote the set of primitive even (resp. odd) Dirichlet characters modulo $N$.
    For $x \in \bR_{> 0}$, let $\lceil x \rceil^{\frp}$ be the smallest prime $\ge x$.
    For $c > 1$, $\delta > 0$, and $y > 0$, define
    \begin{align*}
        P_{\pm}(y, X, c) &:= \frac{\log X}{X} \sum_{\substack{N \in [X, cX] \\ N \text{ prime}}} \sum_{\chi \in \cD_{\pm}(N)} \frac{\chi(\lceil yX \rceil^{\frp})}{\tau(\chi)} \\
        P_{\pm}(y, X, \delta) &:= \frac{\log X}{X^\delta} \sum_{\substack{N \in [X, X + X^\delta] \\ N \text{ prime}}} \sum_{\chi \in \cD_{\pm}(N)} \frac{\chi(\lceil yX \rceil^{\frp})}{\tau(\chi)}.
    \end{align*}
    Then
    \begin{equation}
        \label{eqn:lee_1}
        \lim_{X \to \infty} P_{\pm} (y, X, c) = \begin{cases}
            \int_{1}^{c} \cos\left(\frac{2 \pi y}{x}\right) \dd x & \text{if } + \\
            -i \int_{1}^{c} \sin \left(\frac{2 \pi y}{x}\right) \dd x & \text{if } -
        \end{cases}
    \end{equation}
    and assumming RH, if $\frac{1}{2} < \delta < 1$ we have
    \begin{equation}
        \label{eqn:lee_2}
        \lim_{X \to \infty} P_{\pm} (y, X, \delta) = \begin{cases}
            \cos (2 \pi y) & \text{if } + \\
            -i \sin (2 \pi y) & \text{if } -
        \end{cases}
    \end{equation}
\end{theorem}

The proof is much simpler than the case of modular forms.
The main ingredient of the proof is the following formulas \cite[Lemma 2.6]{lee2025murmurations}: for two distinct primes $p$ and $N$,
\begin{align}
    \sum_{\chi \in \cD_{+}(N)} \frac{\chi(p)}{\tau(\chi)} &= \left(\frac{N-1}{N}\right) \cos \left(\frac{2 \pi p}{N}\right) + \frac{1}{N} \\
    \sum_{\chi \in \cD_{-}(N)} \frac{\chi(p)}{\tau(\chi)} &= -i\left(\frac{N-1}{N}\right) \sin \left(\frac{2 \pi p}{N}\right)
\end{align}
Combined with the prime number theorem (which gives equidistribution results of primes in $[X, cX]$ normalized by $X$), we get \eqref{eqn:lee_1}.
For \eqref{eqn:lee_2}, we need to estimate the number of primes in the short interval $[X, X + X^\delta]$, which is the part that requires RH.

They also proved similar results for real Dirichlet characters, but the proof is more complicated.
Let $\scG$ be the set of odd square-free integers and let $\chi_{d} = \left(\frac{d}{\cdot}\right)$.
For a compactly supported smooth function $\Phi \ge 0$ on $\bR$, define
\begin{equation}
    M_{\Phi}(y, X, \delta) = \frac{\log X}{X^{1 + \delta}} \sum_{\substack{p \in [yX, yX + X^\delta] \\ p \text{ prime}}} \sum_{d \in \scG} \Phi\left(\frac{d}{X}\right) \chi_{8d}(p) \sqrt{p}.
\end{equation}

\begin{theorem}[{Lee--Oliver--Pozdnyakov \cite[Theorem 1.2]{lee2025murmurations}}]
    Fix $y > 0$ and assume $\frac{3}{4} < \delta < 1$.
    Assumming GRH, we have
    \begin{equation}
        M_{\Phi} (y, \delta) := \lim_{X \to \infty} M_{\Phi}(y, X, \delta) = \frac{1}{2} \sum_{\substack{a \ge 1 \\ a \text{ odd}}} \frac{\mu(a)}{a^2} \sum_{m \ge 1} (-1)^{m} \widetilde{\Phi} \left(\frac{m^2}{2 a^2 y}\right),
    \end{equation}
    where
    \begin{equation}
        \widetilde{\Phi}(\xi) = \int_{-\infty}^{\infty} (\cos(2 \pi \xi x) + \sin(2 \pi \xi x)) \Phi(x) \dd x.
    \end{equation}
\end{theorem}
The proof is more involved, which is based on the Polya--Vinogradov inequality 
\[
\left|\sum_{\substack{p \in [yX, yX + X^\delta] \\ p\text{ prime}}} \chi_d(p)\right| \ll (yX)^{\frac{1}{2} + \epsilon}
\]
(for non-principal $\chi_d$ with $\frac{1}{2} < \delta < 1$, which uses GRH \cite{granville2007large})
and a summation formula
\[
\frac{1}{X} \sum_{\substack{d \in \bZ \\ d \text{ odd}}} \left(\sum_{\substack{a^2 | |d| \\ a \le A}} \mu(a)\right) \Phi\left(\frac{d}{X}\right) \left(\frac{d}{p}\right) \sqrt{p} = \frac{1}{2} \left(\frac{2}{p}\right) \sum_{\substack{0 < a \le A \\ (a, 2p) = 1}} \frac{\mu(a)}{a^2} \sum_{k \in \bZ} (-1)^k \left(\frac{k}{p}\right) \widetilde{\Phi} \left(\frac{k X}{2 a^2 p}\right)
\]
which can be proved by using Poisson summation formula.


\subsection{Flying Hecke characters of imaginary quadratic fields}

\cite{wang2025murmurations}

\subsection{Flying modular forms (in weight direction)}

Recall that Zubrilina computed murmuration density for a \emph{fixed weight $k$ and varying level $N$}.
In \cite{bober2023murmurations}, Bober, Booker, Lee\footnote{Min Lee, not Kyu Hwan Lee}, and Lowry-Duda considered the opposite case, where they fix the level $N = 1$ and vary the weight $k$.
In this case, the considered family of Hecke newforms whose \emph{analytic conductor}
\[
\cN(k) := \left(\frac{\exp \psi(k/2)}{2 \pi}\right)^2 \approx \left(\frac{k-1}{4 \pi}\right)^2 + O(1)
\]
are in certain range, where $\psi(x) = \Gamma'(x)/\Gamma(x)$ is the digamma function.

\begin{theorem}[{Bober--Booker--Lee--Lowry-Duda \cite[Theorem 1.1]{bober2023murmurations}}]
    Fix $\epsilon \in (0, \frac{1}{12})$, $\delta \in \{0, 1\}$, and a compact interval $E \subset \bR_{> 0}$ with $|E| > 0$.
    Let $K, H > 0$ with $K^{\frac{5}{6} + \epsilon} < H < K^{1 - \epsilon}$, and let $N = \cN(K)$.
    Then as $K \to \infty$, we have

    \begin{equation}
        \frac{
            \sum_{\substack{p \text{ prime} \\ p / N \in E}}
            \log p
            \sum_{\substack{k \equiv 2 \delta \Mod{4} \\ |k - K| \le H}}
            \sum_{f \in H_k(1)} \lambda_f(p)
        }
        {
            \sum_{\substack{p \text{ prime} \\ p / N \in E}} 
            \log p 
            \sum_{\substack{k \equiv 2 \delta \Mod{4} \\ |k - K| \le H}} 
            \sum_{f \in H_k(1)} 1
        } = \frac{(-1)^\delta}{\sqrt{N}} \left(\frac{\nu(E)}{|E|} + o_{E, \epsilon}(1)\right)
    \end{equation}
    where
    \begin{align}
        \nu(E) &= \frac{1}{\zeta(2)} \sum_{\substack{a, q \in \bZ_{>0} \\ (a, q) = 1 \\ q^2 / a^2 \in E}} \frac{\mu(q)^2}{\varphi(q)^2 \sigma(q)} \left(\frac{a}{q}\right)^{-3} \\
        &= \frac{1}{2} \sum_{t \in \bZ} \prod_{p \nmid t} \frac{p^2 - p - 1}{p^2 - p} \int_E \cos\left(\frac{2 \pi t}{\sqrt{y}}\right) \dd y
    \end{align}
    where the summation $\sum^{*}$ indicates that the terms occuring at the endpoints of $E$ are halved. 
\end{theorem}

The main tool for the proof is the (original) Eichler--Selberg trace formula that does not include Atkin--Lehner operators (e.g. \cite[Theorem 2.1]{child2022twist}).
Then apply class number formula to replace class numbers with the special values of Dirichlet $L$-functions at $s = 1$, which can be estimated under GRH.

\subsection{Flying Maass forms}

Booker, Lee, Lowry-Duda, Seymour-Howell, and Zubrilina computed murmuration densities for weight 0 and level 1 Maass forms \cite{booker2024murmurations}.
They considered a family of Maass forms where the spectral parameter ($R$ with $\lambda = \frac{1}{4} + R^2$) goes ot $\infty$, which is equivalent to the \emph{analytic conductor} $\cN(R)$ going to $\infty$.

\begin{theorem}[{Booker--Lee--Lowry-Duda--Seymour-Howell--Zubrilina \cite[Theorem 1.1]{booker2024murmurations}}]
    Let $E \subset \bR_{>0}$ be a fixed compact interval with $|E| > 0$.
    Let $R, H > 0$ with $R^{\frac{5}{6} + \delta} < H < R^{1 - \delta}$ for some $\delta > 0$ and $N = \cN(R)$.
    Assumming GRH for $L$-functions of Dirichlet characters and Maass forms, as $R \to \infty$ we have 
    \begin{equation}
        \frac{\sum_{\substack{p\text{ prime} \\ p / N \in E}} \log p \sum_{|r(f) - R| \le H} \epsilon(f) a_{f}(p)}{\sum_{\substack{p\text{ prime} \\ p / N \in E}} \log p \sum_{|r(f) - R| \le H} 1} \to \frac{1}{\sqrt{N}|E|}  \sideset{}{^{*}}\sum_{\frac{q^2}{a^2} \in E}  \frac{\mu(q)^2}{\varphi(q)^2 \sigma(q)} \left(\frac{a}{q}\right)^{-3}
    \end{equation}
    where the summation $\sum^{*}$ indicates that the terms occuring at the endpoints of $E$ are halved. 
\end{theorem}

Proof uses an explicit Selberg trace formula due to Str\"ombergsson in his unpublished work \cite{strombergsson}, which requires an analytic test function and cannot be compactly supported, where GRH is needed to control the cutoff error term.
The remaining proof is similar to the weight aspect case of the modular forms \cite{bober2023murmurations}.

\subsection{General formulation}

In fact, all the above works fit into the general framework suggested by Sarnak, in his letter to Sutherland and Zubrilina \cite{sarnak}.

\subsection{Elliptic curves?}

There \emph{are} some works on murmuration density of elliptic curves.

In the last slide of Sutherland's lecture \cite{sutherland} at Tate conference (\emph{The legacy of John Tate, and beyond} at Harvard university), he mentioned an unpublished work with Sawin on murmuration density of elliptic curves, proved by Voronoi summation formula.
He considered it as \emph{a} murmuration theorem, and might not be \emph{the} murmuration theorem since the density formula is too complicated.

\begin{theorem}[Sutherland--Sawin, unpublished]
    Let $\scE(X) := \{y^2 = x^3 + ax + b: a, b \in \bZ, p^4 \mid a \Rightarrow p^6 \nmid b, \max\{4|a|^3, 27 b^2\} \le X\}$ be the set of naive isomorphism classes of elliptic curves over $\bQ$ ordered by height.
    For any smooth function $W : \bR_{>0} \to \bR$ with compact support, the limit
    \begin{equation}
        \lim_{X \to \infty} \frac{1}{|\scE(X)|} \sum_{E \in \scE(X)} \frac{\varepsilon(E)}{N_E} \sum_{n \ge 1} W\left(\frac{n}{N_E}\right) a_n(E)
    \end{equation}
    exists and is equal to
    \begin{equation}
        \int_{0}^{\infty} 2 \pi W(u) \sum_{n \ge 1} \frac{\prod_{p | n} \ell_{p^{\nu_p(n)}}}{\sqrt{n}} \sqrt{u} J_1 (4 \pi \sqrt{un}) \dd u
    \end{equation}
    where $\ell_{p^\nu} = \frac{p^9 - p^8}{p^{10} - 1} \Tr(T_p | S_{\nu + 2}(1))$.
\end{theorem}


\newpage


% --- Bibliography ---

% Start a bibliography with one item.
% Citation example: "\cite{williams}".

\bibliographystyle{acm} % We choose the "plain" reference style
{
% \footnotesize
\small
\bibliography{refs} % Entries are in the refs.bib file
}



% \begin{thebibliography}{1}

% \bibitem{williams}
%    Williams, David.
%    \textit{Probability with Martingales}.
%    Cambridge University Press, 1991.
%    Print.

% % Uncomment the following lines to include a webpage
% % \bibitem{webpage1}
% %   LastName, FirstName. ``Webpage Title''.
% %   WebsiteName, OrganizationName.
% %   Online; accessed Month Date, Year.\\
% %   \texttt{www.URLhere.com}

% \end{thebibliography}

% --- Document ends here ---

\end{document}